\chapter{Related Work}
\label{ch:related}

Where \cref{ch:background} covered the systems this thesis \emph{builds on}, this chapter surveys work that
attacks the \emph{same or similar problems} and positions our contributions against it. We group it into
learned routing and traffic engineering (\cref{sec:rel:routing}), learned congestion control and multipath
scheduling (\cref{sec:rel:transport}), adaptive bitrate and QoE-driven control (\cref{sec:rel:abr}),
goal-conditioned and modulated policies (\cref{sec:rel:gcrl}), and intent-based networking
(\cref{sec:rel:ibn}).

\section{Learned Routing and Traffic Engineering}
\label{sec:rel:routing}

A line of work applies reinforcement learning to routing and traffic engineering. \emph{Learning to
Route}~\cite{Valadarsky2017} studies RL and learned representations for routing, and experience-driven
networking~\cite{Xu2018} uses deep RL for traffic-engineering decisions. These target network-operator
objectives (e.g.\ minimizing congestion globally) from a control-plane vantage point. Our single-path selector differs in
both the actor and the objective: the decision is made by an \emph{endpoint} choosing among
cryptographically authenticated SCION paths, and the objective is the \emph{application's} intent rather than
a global network utility. Most directly related is the DQN SCION path selector of
Bourak~et~al.~\cite{Bourak2025}, which we extend; the differences -- a variable-action scoring architecture
and per-path intent conditioning that removes the single-objective limitation -- are the substance of \cref{ch:single_path}.

The deployed alternative deserves stating alongside the learned one. As \cref{sec:background:pathsel}
describes, a SCION endpoint today reduces its candidate set by policy rather than by measurement: the
path service supplies the segments, a declarative path policy filters them, and a static preference
such as fewest AS hops breaks the remaining ties~\cite{Perrig2017}, with path-aware transport APIs
exposing the same choice to applications. Our \emph{SCION default} baseline
(\cref{sec:p1impl:baselines}) is that rule, with its remaining ties broken on measured latency rather
than statically, which makes it a stronger opponent than the deployed default and not a weaker one.
The distinction our work turns on is that all of these rank
paths on properties that do not change with load, whereas which candidate is best in
\cref{sec:p1design:env} changes between one hour and the next in 8.6\% of contexts.

\section{Learned Congestion Control and Multipath Scheduling}
\label{sec:rel:transport}

Learned transport controllers include Aurora~\cite{Jay2019}, which casts congestion control as RL, and
Orca~\cite{Abbasloo2020}, which combines a learned layer with a classic controller. Mortise~\cite{Shen2026}
is the closest in spirit, and we treat it separately below. For multipath \emph{scheduling}, ReLeS~\cite{Zhang2019} learns a
Multipath QUIC scheduler and Peekaboo~\cite{Wu2020} learns to schedule across heterogeneous paths in dynamic
environments; both optimize a scheduling policy for a given, exogenous load. The multipath system's distinction is that it
does \emph{not} treat the offered load as fixed: it co-optimizes the application's bitrate with the split, so
the controller can relieve congestion it would otherwise create -- a coupling those schedulers, and the
Multipath TCP mechanisms of RFC~8684~\cite{RFC8684}, do not address. For that reason neither is used as
a comparison point in \cref{sec:mp:eval}: neither takes the offered load as a decision, so neither can
be dropped into this system's place without changing the problem it solves.

\paragraph{The fixed rules the field deploys.} Learned schedulers are the minority of what runs. The
default in essentially every Multipath TCP~\cite{RFC8684} and Multipath QUIC~\cite{DeConinck2017}
implementation is minRTT, which fills paths in ascending smoothed RTT, each up to its congestion
window; it is the status quo \cref{ch:multipath} must improve on, and it is the strongest of the five
fixed rules in \cref{tab:p2results}, which the learned pair leads by $+0.139$ to $+0.250$ QoE. Its
known failure mode is head-of-line blocking, when bytes placed on a slower path delay a frame that the
faster paths had already carried, and two refinements address it by predicting that delay before
committing the bytes: BLEST~\cite{Ferlin2016} and ECF~\cite{Lim2017}. Each needs its own send-time
model to be implemented faithfully, and we did not implement them, so the comparison in
\cref{sec:mp:eval} is against the deployed default rather than against the best published fixed rule.

\paragraph{Learned real-time video transport.} The closest published systems to \cref{ch:multipath} are
not multipath schedulers but real-time video controllers. Salsify~\cite{Fouladi2018} couples the video
codec to the transport so that the encoder's rate decision and the packet-scheduling decision are made
together, which is this chapter's premise on a single path. Concerto~\cite{Zhou2019} learns to
coordinate codec and transport for mobile video telephony, and OnRL~\cite{Zhang2020} runs
reinforcement-learned rate control in a deployed video-telephony service at scale, which is the
strongest existing evidence that learned control of this workload survives contact with real networks.
None of them decides how to spread a frame across several paths: their network-side decision is a rate,
not an allocation. That is the delta of \cref{ch:multipath} --- joint rate \emph{and} multipath split ---
and it is also why our comparison point is a WebRTC-style GCC~\cite{Carlucci2016} stack rather than one
of these systems, none of which is defined over a multipath connection.

\subsection{Mortise: intent-driven tuning of a transport mechanism}
The conceptual template for both parts is Mortise~\cite{Shen2026}. Mortise observes that a single
congestion-control algorithm cannot be optimal for every application, and instead \emph{auto-tunes}
the algorithm's parameters to maximize the running application's QoE, using the application's revealed
preferences over throughput, latency, and loss to steer the tuning. The essential idea -- treat
application intent as a first-class input that reshapes a network mechanism rather than a constant
compiled into it -- is exactly what we transplant. The single-path selector applies it to \emph{path selection} (which
path to use), and the multipath system applies it to \emph{joint rate and multipath scheduling} (what
the application sends and how it is spread across paths).

\section{Adaptive Bitrate and QoE-Driven Control}
\label{sec:rel:abr}

Learned adaptive bitrate control, exemplified by Pensieve~\cite{Mao2017}, showed RL beating hand-crafted ABR
heuristics for buffered video streaming, using QoE-shaped rewards and quality metrics such as
VMAF~\cite{VMAF2016}. Our multipath system shares the QoE-driven, learned-bitrate stance but targets \emph{real-time}
video (no playback buffer, deadline-based loss) and, crucially, couples the bitrate decision to a
simultaneous multipath scheduling decision. Where Pensieve adapts rate to a single, exogenous network path,
we adapt rate \emph{and} decide how to spread it across several paths, jointly.

\section{Goal-Conditioned and Modulated Policies}
\label{sec:rel:gcrl}

The machinery for serving many objectives with one policy comes from goal-conditioned RL -- universal value
function approximators~\cite{Schaul2015} and hindsight relabeling~\cite{Andrychowicz2017}, which learn a
value function jointly over states and goals.

The more exact home for this thesis is \emph{multi-objective} RL, since our intent vector is a
preference weighting over a linear scalarization rather than a target state~\cite{Roijers2013}. That
literature already establishes what one might otherwise take to be our claim: a single network
conditioned on the weight vector can represent the family of optima the weight simplex induces and
generalize to weightings it never trained on, whether learned as an envelope value
function~\cite{Yang2019} or under weights that change during execution~\cite{Abels2019}. Our
contribution is therefore not that one conditioned policy can serve many objectives. It is narrower and
sits where that literature does not look: MORL policies emit actions from a fixed, enumerated action
space, in which any conditioning reaching the output layer can change the action, so nothing there asks
\emph{where} in a \emph{comparison-based} scorer the preference must enter. That is the entire content
of \cref{prop:p1cancel,prop:p2cancel}, and the failure it names --- a policy that consumes its objective,
trains on it, and acts identically --- cannot arise in the architectures MORL studies. Hindsight
relabeling does not apply to preference goals at all, since there is no achieved weight vector to
relabel with.

For a comparison-based selector, then, conditioning must
reach the per-alternative terms, because an additive, symmetric signal cancels under the comparison. We
establish this for both forms the comparison takes here --- under an $\arg\max$ over discrete candidates
(\cref{prop:p1cancel}) and under a softmax over continuous allocations (\cref{prop:p2cancel}), where the
cancellation is exact on the whole action rather than only on the winner. We also establish its limit,
which is the more practically useful half: the premise is additivity \emph{on the terms the comparison
consumes}, and a deep encoder that mixes the conditioning with each alternative's features escapes it
without being designed to (\cref{sec:p2eval:conditioning}). The principle is therefore a statement about
which routings carry a guarantee, not about which ones happen to work. This connects a general
representation-learning idea to a concrete failure mode in networked decision-making.

Once the conditioning reaches the per-alternative terms, the mechanism that carries it there appears to
matter little. The obvious richer alternative is feature-wise modulation~\cite{Perez2018}, in which the
objective generates per-feature scales and shifts rather than being concatenated to the features; we
implemented and measured it, found it tied with plain concatenation within the spread between training
seeds, and dropped it because it costs additional parameters for no measurable return. That is the
evidence behind this thesis's claim that the finding is about \emph{where} the intent enters and not
about the machinery that carries it once it is in the right place.

\section{Intent-Based Networking}
\label{sec:rel:ibn}

Finally, ``intent'' also names a body of network-management work concerned with translating high-level
operator goals into configuration, surveyed in RFC~9315~\cite{RFC9315}. That community's ``intent'' is an
operator's declarative policy for a whole network; ours is an \emph{application's} per-flow QoS preference
consumed by a learned endpoint controller. The two are complementary rather than competing -- one could
imagine operator intent bounding the space in which application intent is served -- but they operate at
different layers and timescales, and our techniques address the fine-grained, per-flow end of the spectrum.
